\section{The Laplacian --- Curvature, Diffusion, and Waves}
\label{sec:laplacian}

Applying divergence to the gradient gives
\[
\boxed{\nabla^2 f=\nabla\cdot(\nabla f)
=\frac{\partial^2f}{\partial x^2}+
\frac{\partial^2f}{\partial y^2}+
\frac{\partial^2f}{\partial z^2}.}
\]
The Laplacian compares a value with nearby values in every spatial direction. A positive Laplacian means the point lies below its local average; a negative Laplacian means it lies above it.

\begin{bookfigure}{Positive and negative curvature in one dimension. In several dimensions the Laplacian sums curvature over mutually perpendicular directions.}
\begin{tikzpicture}[x=1.15cm,y=1.05cm]
\draw[physics axis] (-3,0)--(3,0) node[right] {$x$};
\draw[visualcyan!80!black,thick,domain=-2.2:2.2,samples=80] plot(\x,{0.35*\x*\x-1});
\node[font=\small] at (0,-1.45) {$f''>0$};
\begin{scope}[xshift=7cm]
\draw[physics axis] (-3,0)--(3,0) node[right] {$x$};
\draw[visualcyan!80!black,thick,domain=-2.2:2.2,samples=80] plot(\x,{-0.35*\x*\x+1});
\node[font=\small] at (0,1.45) {$f''<0$};
\end{scope}
\end{tikzpicture}
\label{fig:laplacian-curvature}
\end{bookfigure}

\subsection{Harmonic functions}
A function satisfying $\nabla^2 f=0$ is harmonic. Harmonic functions have no local maxima or minima inside a source-free domain unless they are constant. Electrostatic potential in a charge-free region and steady temperature without internal heat sources are harmonic.

\subsection{Diffusion and waves}
The diffusion equation is
\[
\boxed{\frac{\partial u}{\partial t}=D\nabla^2u.}
\]
At a peak, $\nabla^2u<0$, so the peak falls; at a trough, $\nabla^2u>0$, so the trough rises. Diffusion smooths spatial variation. By contrast,
\[
\frac{\partial^2u}{\partial t^2}=c^2\nabla^2u
\]
is a wave equation: the same spatial curvature produces propagation because the time derivative is second order.

\begin{workedexample}
\textbf{Laplacian of a Gaussian.}\par
For $f(x,y)=e^{-(x^2+y^2)}$ and $r^2=x^2+y^2$,
\[
\nabla^2f=4(r^2-1)e^{-r^2}.
\]
It is negative near the central peak and positive outside the unit circle.
\end{workedexample}

\begin{keyidea}[Connection to quantum mechanics]
The nonrelativistic kinetic-energy operator is
\[
\widehat T=-\frac{\hbar^2}{2m}\nabla^2.
\]
The Laplacian therefore measures spatial curvature of the wavefunction and converts that curvature into kinetic energy.
\end{keyidea}
